% -*- coding: utf-8 -*-

\begin{zhixie}

时光荏苒，四年的本科求学生涯即将画上句点。回首这段岁月，心中涌动着深深的感激之情，谨以此文，向所有给予我帮助与支持的人致以最诚挚的谢意。

首先，衷心感谢我的指导教师李重仪教授。李老师在本文的选题确立、研究方向规划以及写作修改的各个阶段，始终给予我悉心的指导与耐心的支持，让我在这段时间里不仅学到了专业知识，更学到了如何认真对待一项工作、如何独立思考与解决问题。这些收获将使我终身受益。
同时，衷心感谢袁捷禹老师在研究过程中给予我的指导与帮助。袁老师在技术实现与实验方案设计方面提供了许多宝贵的建议与细致的指点，帮助我在具体的研究工作中少走弯路。

感谢南开大学计算机学院全体教师在过去四年中对我的悉心培养。各位老师在课堂上传授的专业知识与思维方法，让我打下了较为坚实的理论基础。

特别感谢 4DGaussians 与 SeaSplat 等开源项目的贡献者们。本文的研究工作建立在这些优秀开源代码的基础之上，正是开放共享的学术精神使得相关领域的研究得以快速推进。在此谨向上述项目的开发者致以诚挚的敬意与谢意。

最后，感谢我的家人。在四年的求学历程中，父母始终是我最坚强的后盾，他们无私的爱与默默的付出给了我莫大的力量。无论顺境还是逆境，家人的理解与支持让我能够安心投入学习与研究。能够完成这篇论文，离不开他们多年来的培育与鼓励。

南开大学给予了我四年宝贵的成长空间，“允公允能，日新月异”的校训将激励我在未来的道路上不断进取。谨以此论文，向所有关心和帮助过我的人致以最衷心的感谢。

\end{zhixie}
